\originalchapter{期中练习卷与期末练习卷一}
本章据 OCW 所列 Practice Exam 1, Practice Exam 2, Practice Final Exam 1 及匹配官方解答中文重编并补充推导. 题目中的图均以完整文字约定重述. 分值仅在原件给出时保留. 原题及官方答案中的错误另标为``本稿校正'', 不视为官方结论.
\originalsection{期中练习卷一: Practice Exam 1}
原卷要求闭卷, 禁用电子设备, 答案写在答题册. 公式页的阻尼分界误印为 $\Gamma=\omega_0/2$; 对 $\ddot x+\Gamma\dot x+\omega_0^2x=0$, 特征根为 $-\Gamma/2\pm\sqrt{\Gamma^2/4-\omega_0^2}$, 故本稿校正为 $\Gamma=2\omega_0$.
\begin{exercise}\label{exam:e1.1}
第1题 (30分). 质量 $M$ 静止悬挂在劲度 $k$ 的无质量竖直弹簧下. $t=0$ 时一部分脱落, 仅剩 $\alpha M$, 脱落不给剩余部分冲量. 重力加速度为 $g$, 原平衡位置为 $y=0$, $y$ 向上. (a) 求新平衡位置. (b) 求振动周期. (c) 求 $y(t)$. (d) 求振幅及相位. (e) 求剩余质量的动能、势能随时间的变化.
\end{exercise}
\begin{solution}
官方解答中文重编并补充推导. (a) 原伸长 $\ell_0=Mg/k$. 上移 $y$ 后弹力 $k(\ell_0-y)$, 故 $\alpha M\ddot y=(1-\alpha)Mg-ky$, 新平衡 $y_e=(1-\alpha)Mg/k$.
(b) 令 $q=y-y_e$, 则 $\ddot q+\omega_0^2q=0$, $\omega_0=\sqrt{k/(\alpha M)}$, 周期 $2\pi\sqrt{\alpha M/k}$.
(c),(d) $y(0)=\dot y(0)=0$ 给出 $q=-y_e\cos\omega_0t$, 因而 $y=y_e(1-\cos\omega_0t)$. 以 $q=A\cos(\omega_0t+\phi)$ 定义相位, $A=y_e$, $\phi=\pi$ (模 $2\pi$).
(e) $K=\tfrac12\alpha M y_e^2\omega_0^2\sin^2\omega_0t=\tfrac12ky_e^2\sin^2\omega_0t$. 弹簧实际势能为 $U_s=\tfrac12 k(Mg/k-y)^2$, 重力势能可取 $U_g=\alpha Mgy$. 二者相加、减去新平衡的常数后是 $U=\tfrac12k(y-y_e)^2=\tfrac12ky_e^2\cos^2\omega_0t$. 所以振动能 $K+U=ky_e^2/2$. 本稿校正: 官方把相对新平衡的有效势能称作弹簧势能, 且总能少写因子 $1/2$; 有重力时须区分两者.
\end{solution}
\begin{exercise}\label{exam:e1.2}
第2题 (30分). 水平质量 $m$ 的左侧接阻尼器 $b$, 右侧接弹簧 $k$. 阻尼力按相对速度计算. 方案 A 只移动阻尼器左端, 方案 B 只移动弹簧右端, 移动端位置均为 $s(t)=s_0\cos\omega_dt$, 另一端固定. 对两种方案分别: (a) 写运动方程. (b) 求稳态振幅并画其对频率的变化. (c) 求最大振幅对应频率. (d) 求 $\omega_d\ll\omega_0$ 和 $\omega_d\gg\omega_0$ 的行为.
\end{exercise}
\begin{solution}
官方解答中文重编并补充推导. 取 $x$ 向右, 从所有端点未移动时的平衡量起. 记 $\omega_0^2=k/m$, $\Gamma=b/m>0$.
(a) A: $m\ddot x=-kx-b(\dot x-\dot s)$; B: $m\ddot x=-k(x-s)-b\dot x$.
(b) 以 $x=\Rea(C\e^{-\ii\omega_dt})$ 表示, 分母 $D=\omega_0^2-\omega_d^2-\ii\Gamma\omega_d$, 有 $C_A=-\ii\Gamma\omega_ds_0/D$, $C_B=\omega_0^2s_0/D$. 因而
\[
 A_A=\frac{\Gamma\omega_ds_0}{\sqrt{(\omega_0^2-\omega_d^2)^2+\Gamma^2\omega_d^2}},\qquad
 A_B=\frac{\omega_0^2s_0}{\sqrt{(\omega_0^2-\omega_d^2)^2+\Gamma^2\omega_d^2}}.
\]
(c) 对振幅平方求导即可避免根号. A 的极值条件是 $\omega_0^4-\omega_d^4=0$, 唯一正峰在 $\omega_d=\omega_0$, 高度 $s_0$. B 的分母平方对 $u=\omega_d^2$ 为 $(\omega_0^2-u)^2+\Gamma^2u$, 极小点 $u=\omega_0^2-\Gamma^2/2$. 若 $\Gamma<\sqrt2\omega_0$, 峰在此正频率; 否则最大值在零频率, 曲线单调下降.
(d) A 从零线性起升, $A_A\sim\Gamma s_0\omega_d/\omega_0^2$, 过唯一峰后以 $\Gamma s_0/\omega_d$ 衰减. B 从 $s_0$ 起, 弱阻尼时先上升再下降, 强阻尼时直接下降, 高频 $A_B\sim\omega_0^2s_0/\omega_d^2$. 这些起点、峰和尾部给出所需草图. 本稿校正: 官方峰频率漏写 $\Gamma$ 的平方并混用 $1/2$, 低频 B 振幅误写为 $\omega_0^2s_0$, 且遗漏强阻尼情形.
\end{solution}
\begin{exercise}\label{exam:e1.3}
第3题 (40分). 对以下五个系统定义坐标, 写小振幅运动方程, 化为 $(\mathsf A-\omega^2\mathsf I)\boldsymbol C=0$. 机械系统 $\mathsf A=\mathsf M^{-1}\mathsf K$, 电路用等效矩阵. 不要求求解本征值.
(A) 水平序列为固定端--$k$--$m$--$k$--$3m$--$k$--$m$--$k$--固定端.
(B) 三个质量 $m$ 各由长 $L$ 的摆线悬挂, 只保留水平运动; 左中间弹簧 $k$, 中右间弹簧 $3k$.
(C) 四个质量 $m$ 沿圆环运动, 忽略环曲率. 上、下两条相邻连边弹簧为 $k$, 左、右两条为 $3k$.
(D) 上支路三个电感 $L$ 串联, 四个节点各经电容 $C$ 接共同下支路; 用三个电感电流为坐标.
(E) 顶端固定的三连摆, 三段无质量摆杆各长 $L$, 各接一个质量 $m$; 只考虑小角度.
\end{exercise}
\begin{solution}
官方解答中文重编并补充推导. 所有坐标满足 $\ddot{\boldsymbol x}+\mathsf A\boldsymbol x=0$, 代入 $\boldsymbol x=\boldsymbol C\e^{-\ii\omega t}$ 即得题定矩阵方程.
(A) 左至右 $x_j$ 向右. 弹簧伸长为 $x_1,x_2-x_1,x_3-x_2,-x_3$. 由各力的差得
\[
\mathsf M=\operatorname{diag}(m,3m,m),\quad
\mathsf K=k\begin{pmatrix}2&-1&0\\-1&2&-1\\0&-1&2\end{pmatrix},\quad
\mathsf A=\frac{k}{m}\begin{pmatrix}2&-1&0\\-1/3&2/3&-1/3\\0&-1&2\end{pmatrix}.
\]
(B) 左至右 $x_j$ 向右, 单摆回复力 $-mgx_j/L$, 弹簧力由相邻位移差给出, 所以
\[
\mathsf A=\frac gL\mathsf I+\frac km\begin{pmatrix}1&-1&0\\-1&4&-3\\0&-3&3\end{pmatrix}.
\]
(C) 自右上质量起顺时针编号, $x_j$ 沿顺时针切向. 每点连接的刚度和为 $4k$, 相邻耦合为负刚度, 因而
\[
\mathsf A=\frac km\begin{pmatrix}4&-3&0&-1\\-3&4&-1&0\\0&-1&4&-3\\-1&0&-3&4\end{pmatrix}.
\]
(D) 电流 $I_j$ 向右, 四个电容上板电荷 $Q_j$. 节点方程为 $\dot Q_1=-I_1$, $\dot Q_2=I_1-I_2$, $\dot Q_3=I_2-I_3$, $\dot Q_4=I_3$. 电感方程 $L\dot I_j=(Q_j-Q_{j+1})/C$, 微分后
\[
\ddot{\boldsymbol I}+\frac1{LC}\begin{pmatrix}2&-1&0\\-1&2&-1\\0&-1&2\end{pmatrix}\boldsymbol I=0.
\]
(E) $x_j$ 为自顶至底三个质量的水平位移, $x_0=0$. 静张力依次为 $3mg,2mg,mg$. 每杆倾角 $(x_j-x_{j-1})/L$, 每质点受上下杆水平张力差. 所以
\[
\mathsf A=\frac gL\begin{pmatrix}5&-2&0\\-2&3&-1\\0&-1&1\end{pmatrix}.
\]
本稿校正: 官方 A 的刚度矩阵前因子误记为 $\omega_0^2$ 而应为 $k$; D 的第三个回路电流符号在中间式不一致, 以上统一向右并直接由节点关系推导.
\end{solution}
\originalsection{期中练习卷二: Practice Exam 2}
本卷原题未给逐题分值.
\begin{exercise}\label{exam:e2.1}
第1题. 等离子体的一维电荷密度扰动满足 $\rho_{tt}=c^2\rho_{xx}-\omega_p^2\rho$, $c$ 为光速, $\omega_p$ 为固定等离子体频率. (a) 对 $a\sin(kx-\omega t)$ 求色散关系. (b) 求 $k(\omega)$, 画波数、相速度、群速度曲线. (c) 验证 $a\cos\omega t\e^{-\kappa x}$ 也是解, 求并画 $\kappa(\omega)$.
\end{exercise}
\begin{solution}
官方解答中文重编并补充推导. (a) 代入得 $-\omega^2=-c^2k^2-\omega_p^2$, 即 $\omega=\sqrt{c^2k^2+\omega_p^2}$.
(b) 正向行波 $k=\sqrt{\omega^2-\omega_p^2}/c$ ($\omega>\omega_p$). 曲线从 $(\omega_p,0)$ 起, 切线初始竖直, 大频率渐近 $\omega/c$. $v_p=\omega/k=c/\sqrt{1-(\omega_p/\omega)^2}$ 从无穷降至 $c$; $v_g=\dd\omega/\dd k=c\sqrt{1-(\omega_p/\omega)^2}$ 从零升至 $c$, 且 $v_pv_g=c^2$.
(c) 代入得 $-\omega^2=c^2\kappa^2-\omega_p^2$, 因而 $\kappa=\sqrt{\omega_p^2-\omega^2}/c$, $0\le\omega<\omega_p$. 图为第一象限椭圆弧, 从 $(0,\omega_p/c)$ 降至 $(\omega_p,0)$. 这是倏逝场而非具有实波数的传播波.
\end{solution}
\begin{exercise}\label{exam:e2.2}
第2题. 很长的弦在 $x=0$ 接无质量小环, 小环可沿竖直杆无摩擦运动. 弦线密度 $\mu$, 张力 $T$. 外力 $F_0\sin\omega t$ 竖直作用于小环, 产生沿正 $x$ 的稳态波, 不考虑返回的反射波. (a) 写端点边界条件. (b) 求产生的波频率、波长. (c) 求振幅.
\end{exercise}
\begin{solution}
官方解答中文重编并补充推导. (a) 小环质量为零, 小斜率下向上张力为 $Ty_x(0,t)$, 所以 $Ty_x(0,t)+F_0\sin\omega t=0$.
(b) 波速 $v=\sqrt{T/\mu}$, 稳态角频率与驱动同为 $\omega$, 普通频率 $\omega/(2\pi)$, 波长 $\lambda=2\pi v/\omega$.
(c) 正向波写 $y=A\cos[\omega(t-x/v)+\phi]$, 边界给 $TA\omega\sin(\omega t+\phi)/v=-F_0\sin\omega t$. 可取 $\phi=\pi$, $A=F_0v/(T\omega)=F_0/(\omega\sqrt{T\mu})$. 该式同时满足波动方程及无入射波条件.
\end{solution}
\begin{exercise}\label{exam:e2.3}
第3题. 长 $L$ 的弦左端 $x=0$ 接无质量环, 可沿垂直弦的无摩擦杆滑动; 右端 $x=L$ 固定. 张力 $T$, 线密度 $\mu$, 忽略重力. 初位移在 $[L/4-D/2,L/4+D/2]$ 上为 $H$, 其余为零, $0<D\ll L$, 初速度处处为零. (a) 求各模态的函数、波数、频率、周期. (b) 求初位移展开系数 $A_m$, 算完积分并简化. (c) 哪些系数为零? (d) 写完整时间级数. (e) 在最低模态周期的一半时求弦形状, 可不用傅里叶展开.
\end{exercise}
\begin{solution}
官方解答中文重编并补充推导. (a) 左端力平衡 $y_x(0,t)=0$, 右端 $y(L,t)=0$. 因此空间函数 $\cos k_mx$, $k_m=(m-\tfrac12)\pi/L$, $m=1,2,\ldots$, $v=\sqrt{T/\mu}$, $\omega_m=vk_m$, $\tau_m=2\pi/(vk_m)$.
(b) 正交范数为 $L/2$, 故
\[
A_m=\frac{2H}{L}\int_{L/4-D/2}^{L/4+D/2}\cos k_mx\dd x
=\frac{4H}{Lk_m}\cos(k_mL/4)\sin(k_mD/2).
\]
(c) $\cos[(m-1/2)\pi/4]$ 不可能为零, 因其零点要求 $m-1/2=2+4j$, 与整数 $m$ 矛盾. 仅当 $(m-1/2)D/(2L)$ 为正整数时 $A_m=0$, 是否存在取决于 $D/L$.
(d) $y=\sum_{m\ge1}A_m\cos(k_mx)\cos(\omega_mt)$; 无正弦时间项, 因初速度为零. 跳跃处级数取左右平均值.
(e) $\tau_1=4L/v$, 在 $t=2L/v$ 所有时间因子均为 $\cos[(2m-1)\pi]=-1$, 故弦仍在原矩形支集上, 高度变为 $-H$. 等价地两个初始半高脉冲各绕行, 固定端反射变号, 自由端反射不变号.
\end{solution}
\begin{exercise}\label{exam:e2.4}
第4题. 左弦 $x<0$ 向左无限延伸, $\mu_1=\mu,T_1=T$; 右弦 $0<x<L$ 在 $L$ 固定, $\mu_2=\mu/2,T_2=T/2$. 两弦由无质量环在 $x=0$ 相接, 环沿竖直无摩擦杆移动. 初始三角脉冲向右运动, 宽 $L/4$; 在 $t=0$ 前沿为 $-7L/8$, 后沿为 $-9L/8$. 原图是后沿零、向前沿线性升高后骤降的斜坡三角形. (a) 写两侧波方程和连接条件. (b) 在 $0,L$ 是否反射, 是否变号? (c) 给入射波 $f(x/v_1-t)$, 求透射和反射函数, 只考虑 $t\le2L\sqrt{\mu/T}$. (d) 画 $t=L\sqrt{\mu/T}$ 及 $2L\sqrt{\mu/T}$ 时形状.
\end{exercise}
\begin{solution}
官方解答中文重编并补充推导. 记 $v=\sqrt{T/\mu}$, 两侧波速相同. (a) 两侧均有 $y_{tt}=v^2y_{xx}$, 连接处 $y_L(0,t)=y_R(0,t)$, $Ty_{L,x}(0,t)=(T/2)y_{R,x}(0,t)$, 固定端 $y_R(L,t)=0$.
(b) 阻抗 $Z_1=\sqrt{T\mu}$, $Z_2=Z_1/2$, 因而 $r=(Z_1-Z_2)/(Z_1+Z_2)=1/3$, 同号反射; 固定端 $r_L=-1$, 反号.
(c) 令透射系数 $q$, 连续条件 $1+r=q$, 力条件 $1-r=q/2$, 解得 $q=4/3$. 在给定时间内固定端反射尚未返回接点, 故
\[
y_L=f(x/v-t)+\tfrac13f(-x/v-t),\quad
y_R=\tfrac43\{f(x/v-t)-f[(2L-x)/v-t]\}.
\]
(d) 以下用精确分段函数代替草图. 令 $\ell=L/8$, 初始峰高 $h$ (原题未指定), $P(u)=h(u+9L/8)/(L/4)$ 在 $-9L/8\le u<-7L/8$, 其余为零; 则 $f(s)=P(vs)$. 在 $t=L/v$, 左侧 $[-\ell,0]$ 形状为 $P(x-L)+P(-x-L)/3$, 即从 $h/3$ 线性升至接点 $2h/3$, 在 $-\ell$ 向左骤降为零; 右侧 $[0,\ell]$ 为 $4P(x-L)/3$, 从 $2h/3$ 升至 $4h/3$ 后骤降. 在 $t=2L/v$, 左侧 $[-9L/8,-7L/8]$ 为 $P(-x-2L)/3$, 是峰高 $h/3$ 的反向斜坡; 右侧 $[L-\ell,L]$ 为 $\tfrac43[P(x-2L)-P(-x)]$, 从 $-4h/3$ 线性升至固定端零, 其余为零. 端点跃变按原脉冲的单侧极限理解.
\end{solution}
\begin{exercise}\label{exam:e2.5}
第5题. 长 $L$, 总质量 $M$, 张力 $T$ 的弦初始两端固定. (a) 求最低模态波长和角频率. (b) 求第 $n$ 模态波长和角频率. 将弦拉成关于 $L/2$ 对称的窄等边三角形, 顶点高 $A$, 其余弦平直; 原图竖直尺度夸张, 实际采用小斜率弦模型. 初形为 $f(x)$, $B_n=(2/L)\int_0^Lf(x)\sin(n\pi x/L)\dd x$. (c) 不积分, 判断哪些 $B_n$ 为零. 从静止释放. (d) 求首次全弦平直及首次恢复初形的时间. (e) 画释放后 $\tfrac14\sqrt{ML/T}$ 的弦形. (f) 改以首次全弦平直为 $t=0$, 已知初形的 $B_5$, 求 $y_5(x,t)$, 不计算 $B_5$. 一位朋友求出系数并叠加最低20个允许模态, 很好拟合初形. (g) 将该级数延拓至 $0\le x\le3L$, 画图. 现在右端改接无摩擦竖直滑杆上的无质量环, 张力不变. (h),(i) 在新配置下重答(a),(b). (j) 初形仍为上述三角形, 将新的模态级数延拓到 $3L$ 并画图.
\end{exercise}
\begin{solution}
官方解答中文重编并补充推导. 记 $v=\sqrt{TL/M}$. (a),(b) 两端节点要求 $k_n=n\pi/L$, 所以 $\lambda_n=2L/n$, $\omega_n=n\pi\sqrt{T/(ML)}$, 最低模态取 $n=1$.
(c) $f(L-x)=f(x)$, 而 $\sin[n\pi(L-x)/L]=(-1)^{n+1}\sin(n\pi x/L)$. 奇偶对称性使所有偶数 $n$ 的积分为零. 不能仅凭此宣称全部奇数系数非零; 某些取决于三角形宽度的额外零点仍可能存在.
(d) $y=\sum B_n\sin(n\pi x/L)\cos(n\pi vt/L)$, 仅有奇数模态. 在 $t=L/(2v)=\tfrac12\sqrt{ML/T}$ 所有时间因子为零. 因 $B_1>0$, 再早不可能全弦为零. 首次恢复为 $t=2L/v=2\sqrt{ML/T}$, 也由非零基频限制.
(e) 将初形以两端奇反射作 $2L$ 周期延拓 $F$, 达朗贝尔解为 $y=[F(x-vt)+F(x+vt)]/2$. 在 $t=L/(4v)$, 原三角形分成两个向左右移动的半高三角形, 中心在 $L/4,3L/4$, 峰高 $A/2$, 宽度仍与初三角形相同 (窄脉冲尚未碰端点). 这给出精确草图.
(f) 实际释放时间为 $t+L/(2v)$, 故
\[
y_5=B_5\sin(5\pi x/L)\cos[5\pi vt/L+5\pi/2]
=-B_5\sin(5\pi x/L)\sin[5\pi t\sqrt{T/(ML)}].
\]
(g) 延拓在 $[0,L]$ 为正三角形, 在 $[L,2L]$ 为中心 $3L/2$ 的负三角形, 在 $[2L,3L]$ 为中心 $5L/2$ 的正三角形. 准确关系是 $F(x+L)=-F(x)$, $F(x+2L)=F(x)$.
(h),(i) 新条件 $y(0)=0,y_x(L)=0$ 给 $k_n=(2n-1)\pi/(2L)$, $\lambda_n=4L/(2n-1)$, $\omega_n=(2n-1)\pi\sqrt{T/(ML)}/2$. 最低模态取 $n=1$.
(j) 右端偶反射给 $F(2L-x)=F(x)$, 再有 $F(x+2L)=-F(x)$. 因此 $[0,L]$, $[L,2L]$, $[2L,3L]$ 的三角形依次正、正、负, 中心仍在各区间中点. 有限20项只近似这些折线, 但反射关系精确.
本稿校正: 官方把首次平直时间取成 $L/(4v)$, 造成(e),(f)错误; 另在(g)写 $F(x+2L)=-F(x)$, 与其同一行及图示矛盾.
\end{solution}
\begin{exercise}\label{exam:e2.6}
第6题. 理想导体平面位于 $z=0,L$, 其间真空有驻电磁波, 电场始终沿 $y$, 磁场沿 $x$. (a) 用导体电边界写第 $n$ 模态电场. (b) 用麦克斯韦方程或两个行波叠加求磁场. (c) 求坡印廷矢量、电磁能量密度及各自时间平均. (d) 第2模态在 $t=0$ 时, 画电场、磁场、电能密度、磁能密度沿 $z$ 的变化.
\end{exercise}
\begin{solution}
官方解答中文重编并补充推导. (a) 导体表面切向电场为零. 选时间原点使电场最大. 写成
\[E_y=E_0\sin(k_nz)\cos\omega_nt,\qquad k_n=n\pi/L,\quad\omega_n=ck_n.\]
(b) 法拉第方程 $\partial_tB_x=\partial_zE_y$ 给 $B_x=(E_0/c)\cos(k_nz)\sin\omega_nt$, 不含额外静磁场.
(c) $\hat y\times\hat x=-\hat z$, 所以
\[
\boldsymbol S=-\frac{E_0^2}{\mu_0c}\sin k_nz\cos k_nz\cos\omega_nt\sin\omega_nt\,\hat z,
\]
\[
u_E=\frac{\epsilon_0E_0^2}{2}\sin^2k_nz\cos^2\omega_nt,\qquad
u_B=\frac{\epsilon_0E_0^2}{2}\cos^2k_nz\sin^2\omega_nt.
\]
时间平均为 $\langle\boldsymbol S\rangle=0$, $\langle u_E\rangle=\epsilon_0E_0^2\sin^2k_nz/4$, $\langle u_B\rangle=\epsilon_0E_0^2\cos^2k_nz/4$, 总平均均匀为 $\epsilon_0E_0^2/4$.
(d) $n=2,t=0$ 时 $E_y=E_0\sin(2\pi z/L)$: 在 $0,L/2,L$ 为零, $L/4$ 正峰, $3L/4$ 负峰. $u_E$ 在两峰处均为 $\epsilon_0E_0^2/2$, 三节点为零; $B_x=u_B=0$ 处处. 本稿校正: 官方坡印廷式漏写 $1/c$, 磁能平均下标误写电能.
\end{solution}
\originalsection{期末练习卷一: Practice Final Exam 1}
此名称来自 OCW 门户的 Practice Final Exam 1. 原 PDF 内封面实际写``Exam 1'', 未注明实际考试年份, 故不把它称作2016年实际期末考试. 原卷要求闭卷, 全做六题, 展示过程, 禁用计算器、书本、电脑和手机. 封面总分表的第5题为18分, 题页标题却为20分; 该题小问分值之和为18分, 以下同时保留这一差异.
\begin{exercise}\label{exam:f1.1}
第1题 (16分, 每问4分). 无限长理想导体矩形波导沿 $x$ 延伸, 截面 $0<y<a,0<z<b$, $a<b$. 给定 TE 模态
\[
\boldsymbol E=E_0\cos(k_yy+\phi_y)\cos(k_xx-\omega t)\hat z.
\]
(a) 求满足边界的 $k_y,\phi_y$. (b) 写色散关系. (c) 求该模态能传播的最低频率. (d) 求相应磁场.
\end{exercise}
\begin{solution}
官方解答中文重编并补充推导. (a) 在 $y=0,a$, 电场沿壁切向, 必须为零, 可取 $\phi_y=-\pi/2$, $k_y=n\pi/a$, $n\ge1$. 在 $z$ 壁电场为法向, 不要求为零. 本稿校正: 官方(a)把 $y=a$ 误写为 $y=b$, 其 $k_y=n\pi/a$ 结果正确. 于是 $E_z=E_0\sin k_yy\cos\psi$, $\psi=k_xx-\omega t$.
(b) 真空波方程代入给 $\omega^2=c^2(k_x^2+k_y^2)$.
(c) 第 $n$ 横向模态截止角频率为 $\omega_{c,n}=n\pi c/a$, 普通频率 $nc/(2a)$. 非零纵向传播要求 $\omega>\omega_{c,n}$; 等号是 $k_x=0$ 的截止点. 题给这族模态最低截止为 $\pi c/a$; 这不是所有可能偏振的矩形波导模态的全局最低截止.
(d) $\nabla\times\boldsymbol E=(\partial_yE_z)\hat x-(\partial_xE_z)\hat y=-\partial_t\boldsymbol B$, 积分得到
\[
\boldsymbol B=\frac{E_0}{\omega}\left[k_y\cos(k_yy)\sin\psi\,\hat x-k_x\sin(k_yy)\cos\psi\,\hat y\right].
\]
取无静态附加磁场. 检查 $\nabla\cdot\boldsymbol B=0$, 导体壁的法向磁场为零, 安培方程则重新给出(b)的色散式.
\end{solution}
\begin{exercise}\label{exam:f1.2}
第2题 (16分). 下方质量 $2m$ 以劲度 $2k$ 的水平弹簧接左固定墙, 上方质量 $m$ 以劲度 $k$ 的水平弹簧接右固定墙. 两质点在平衡时具有相同 $x$ 坐标, 弹簧均未压缩; 其间竖直弹性带长 $L$, 张力 $T=2kL$. 只能沿 $x$ 运动, 无摩擦、无重力, $|x_1|,|x_2|\ll L$, 可认为张力恒定. 下方位移为 $x_1$, 上方为 $x_2$, 均向右为正. (a,5分) 写耦合运动方程. (b,7分) 求模态频率. (c,4分) 画两模态, 明确方向和相对幅度.
\end{exercise}
\begin{solution}
官方解答中文重编并补充推导. (a) 带斜角近似 $(x_2-x_1)/L$, 下方水平带力为 $2k(x_2-x_1)$, 上方相反. 因而
\[
2m\ddot x_1=-4kx_1+2kx_2,\qquad m\ddot x_2=2kx_1-3kx_2.
\]
(b) $\mathsf A=(k/m)\begin{pmatrix}2&-1\\-2&3\end{pmatrix}$. 令 $u=m\omega^2/k$, 特征方程 $(2-u)(3-u)-2=(u-1)(u-4)=0$, 得 $\omega_- =\sqrt{k/m}$, $\omega_+=2\sqrt{k/m}$.
(c) 低频模态 $(C_1,C_2)=(1,1)$, 两质点同向同幅, 带不倾斜. 高频模态 $(1,-2)$, 下方向右一个单位时上方向左两个单位; 过半周期方向同时反转. 这组箭头及比例即为所需模态图.
\end{solution}
\begin{exercise}\label{exam:f1.3}
第3题 (16分). (a,5分) 光纤芯是折射率 $n_f$ 的实心圆柱, 外包折射率 $n_c$ 的圆柱壳. 光从折射率 $n_a$ 的空气入射端面, 入射角 $\theta$ 相对于光纤轴线 (端面法线) 测量. 求使光始终留在芯内的最大 $\theta$. (b,4分) 真空中非偏振光从液体表面反射, 反射光到达水平距离25 cm、竖直高度20 cm的屏幕点, 完全偏振, 求液体折射率 $n$. (c,7分) 色散 $\omega^2=\omega_0^2+A^2k^2$, 常数 $\omega_0,A$ 为实数. (i) 求可传播频率范围. (ii) 求相速度、群速度并画其随 $\omega$ 的标注草图.
\end{exercise}
\begin{solution}
官方解答中文重编并补充推导. (a) 假设 $n_f>n_c$. 芯内光线与轴夹角 $\alpha$, 芯壳界面入射角为 $\pi/2-\alpha$. 全反射要求 $\cos\alpha\ge n_c/n_f$, 故 $\sin\alpha\le\sqrt{n_f^2-n_c^2}/n_f$. 端面斯涅尔定律 $n_a\sin\theta=n_f\sin\alpha$ 给
\[
\theta_{\max}=\arcsin\!\left(\min\{1,\sqrt{n_f^2-n_c^2}/n_a\}\right).
\]
等号是临界极限; 若 $n_f\le n_c$, 无这种全反射导光机制. 原官方公式未说明接受角饱和至 $90^\circ$ 的情形.
(b) 完全偏振说明入射为布儒斯特角, 相对于竖直法线 $\tan\theta_B=25/20$. 真空入射满足 $n=\tan\theta_B=5/4$.
(c) 设 $A\ne0$, 正角频率和正向波数. $k=\sqrt{\omega^2-\omega_0^2}/|A|$, 传播要求 $\omega>|\omega_0|$, 等号为截止点. $v_p=|A|/\sqrt{1-\omega_0^2/\omega^2}$, $v_g=|A|\sqrt{1-\omega_0^2/\omega^2}$. 相速从无穷单调下降到 $|A|$, 群速从零单调上升到 $|A|$; 横轴起点为 $|\omega_0|$, 两曲线有共同水平渐近线. 若 $\omega_0=0$, 两者恒为 $|A|$. 本稿校正: 官方反解波数的一行仍写了含 $k$ 的原式, 应为上述 $k(\omega)$; 绝对值处理保留原题``实常数''的范围.
\end{solution}
\begin{exercise}\label{exam:f1.4}
第4题 (16分). 单色光照射 $N$ 条等宽狭缝, 缝宽 $D$, 相邻缝中心距 $d$, 屏距远大于装置尺寸. 原图横轴 $\sin\theta$, 纵轴强度: 相邻主极大之间有三个小极大, 中央衍射包络内可见第0、$\pm1,\pm2,\pm3$ 级主极大, 第 $\pm4$ 级在第一个包络零点处缺失. (a,6分) 根据图判断 $N$, 解释理由. (b,6分) 求 $d/D$, 解释理由. (c,4分) 将 $D$ 缩小至趋近零, 同时增大入射强度使中央峰不变; 在原图虚线上画新强度曲线.
\end{exercise}
\begin{solution}
官方解答中文重编并补充推导. (a) $N$ 缝干涉因子 $[\sin N\beta/\sin\beta]^2$, $\beta=\pi d\sin\theta/\lambda$. 两主峰间有 $N-1$ 个零点及 $N-2$ 个次峰, 故 $N-2=3$, $N=5$.
(b) 干涉主峰 $d\sin\theta=m\lambda$, 包络首零点 $D\sin\theta=\lambda$. 第4级缺失给 $4\lambda/d=\lambda/D$, 因此 $d/D=4$.
(c) 用归一化中央强度 $I_c$ 表示精确图形:
\[
\frac{I_{\rm old}}{I_c}=\left(\frac{\sin\alpha}{\alpha}\right)^2\left(\frac{\sin5\beta}{5\sin\beta}\right)^2,
\quad \alpha=\beta/4;\qquad
\frac{I_{\rm new}}{I_c}=\left(\frac{\sin5\beta}{5\sin\beta}\right)^2.
\]
可去奇点按连续极限取值. 新图所有主峰均高 $I_c$, 原先缺失的第4级重新出现; 干涉因子主级位置和每两主级间三次峰的结构保留; 有限缝宽时实际总强度极大可因包络而略偏移, 宽衍射包络不再压低远处峰.
\end{solution}
\begin{exercise}\label{exam:f1.5}
第5题 (封面18分, 题页写20分). 长 $2L$ 的弦两端固定, 线密度 $\mu$, 张力 $T$. $a\ll L$, 初位移为零, 初速度为
\[
u(x)=\dot y(x,0)=\begin{cases}v_0,&L-a\le x<L,\\-v_0,&L\le x<L+a,\\0,&\text{其余位置}.\end{cases}
\]
(a,3分) 不论是否被激发, 画前三个模态. (b,10分) 求第 $n$ 模态的振幅, 哪个最低模态未激发? (c,5分) 画 $t=(L/2)\sqrt{\mu/T}$ 时弦位移.
\end{exercise}
\begin{solution}
官方解答中文重编并补充推导. 记 $v=\sqrt{T/\mu}$. (a) 空间模态 $\sin(n\pi x/(2L))$. 前三模态分别有1、2、3个半波, 内节点分别无、$L$、$2L/3,4L/3$.
(b) $k_n=n\pi/(2L)$, $\omega_n=vk_n$. 初位移为零, 可写 $y=\sum A_n\sin k_nx\sin\omega_nt$, 其中 $A_n$ 为有符号系数, 物理振幅 $|A_n|$. 用范数 $L$ 得
\begin{align*}
A_n\omega_n&=\frac{v_0}{L}\left(\int_{L-a}^{L}\sin k_nx\dd x-\int_L^{L+a}\sin k_nx\dd x\right)\\
&=-\frac{2v_0}{Lk_n}\cos(k_nL)[1-\cos(k_na)],\\
A_n&=-\frac{8Lv_0}{vn^2\pi^2}\cos(n\pi/2)\left[1-\cos\frac{n\pi a}{2L}\right].
\end{align*}
所有奇数 $n$ 未激发, 最低为 $n=1$. 对特定 $a/L$ 还可能有偶数模态的额外零点. 这也符合初速度关于中心反对称, 只激发反对称空间模态.
(c) 达朗贝尔初速度解 $y(x,t)=(2v)^{-1}\int_{x-vt}^{x+vt}u(s)\dd s$, 本时刻脉冲未抵达端点, 无需反射. 令
\[
G(s)=\begin{cases}v_0(s-L+a),&L-a\le s\le L,\\v_0(L+a-s),&L\le s\le L+a,\\0,&\text{其余位置},\end{cases}
\]
则 $G'=u$, $y(x,L/(2v))=[G(x+L/2)-G(x-L/2)]/(2v)$. 左侧是中心 $L/2$、支集 $[L/2-a,L/2+a]$ 的正三角形, 右侧中心 $3L/2$ 为同宽负三角形; 峰值分别 $\pm av_0/(2v)$. 本稿校正: 官方文字峰高 $av_0/v$ 漏掉达朗贝尔分裂的 $1/2$ 因子.
\end{solution}
\begin{exercise}\label{exam:f1.6}
第6题 (18分). 质量 $M$、电荷 $+Q$ 的小粒子接在沿 $x$ 轴的劲度 $k$ 弹簧末端, 平衡位于原点. 将其沿正 $x$ 移至 $A$ 后在 $t=0$ 静止释放, 粒子尺寸远小于 $A$, 辐射阻尼很弱. (a,4分) 求 $x$--$z$ 平面任意方向、距离 $R\gg A$ 的辐射电场. (b,4分) 求时间平均总辐射功率. (c,6分) 假定功率随时间变化可忽略, 粗估振幅衰减到 $A/\e$ 的时间, 此假设是否现实? (d,4分) 用 $\dd A/\dd t=(\dd A/\dd E)(\dd E/\dd t)$ 和逐周期平均功率改进估计.
\end{exercise}
\begin{solution}
官方解答中文重编并补充推导. 在非相对论、远区电偶极近似下, $\omega_0=\sqrt{k/M}$, $x=A\cos\omega_0t$, $\boldsymbol a=-A\omega_0^2\cos\omega_0t\,\hat x$. (a) 观察方向 $\hat n=\sin\theta\,\hat x+\cos\theta\,\hat z$, $\theta$ 从正 $z$ 轴量起, 推迟时间 $t_R=t-R/c$. 辐射场
\[
\boldsymbol E_{\rm rad}=\frac{Q}{4\pi\epsilon_0c^2R}\hat n\times(\hat n\times\boldsymbol a(t_R))
=\frac{QA\omega_0^2}{4\pi\epsilon_0c^2R}\cos\omega_0t_R
(\cos^2\theta\,\hat x-\sin\theta\cos\theta\,\hat z).
\]
沿振动轴辐射为零, 垂直方向最大. $\boldsymbol B=\hat n\times\boldsymbol E/c$.
(b) 拉莫尔公式 $P=Q^2a^2/(6\pi\epsilon_0c^3)$, 平均 $\cos^2$ 为 $1/2$, 因而 $\langle P\rangle=Q^2A^2\omega_0^4/(12\pi\epsilon_0c^3)$.
(c) 初能 $E_i=M\omega_0^2A^2/2$, 目标能 $E_i/\e^2$. 固定用初始功率给
\[
t_{\rm rough}=\frac{E_i(1-\e^{-2})}{\langle P\rangle_i}
=\frac{6\pi\epsilon_0Mc^3}{Q^2\omega_0^2}(1-\e^{-2}).
\]
系数 $1-\e^{-2}\simeq0.865$. 振幅减小时功率按振幅平方减小, 所以恒功率假设不现实, 此估计偏短.
(d) 令慢变振幅 $\mathcal A(t)$, 能量 $E=M\omega_0^2\mathcal A^2/2$, 则
\[
\dot E=-\frac{Q^2\omega_0^4\mathcal A^2}{12\pi\epsilon_0c^3},\qquad
\dot{\mathcal A}=-\frac{Q^2\omega_0^2}{12\pi\epsilon_0Mc^3}\mathcal A.
\]
故 $\mathcal A=A\e^{-t/\tau_A}$, $\tau_A=12\pi\epsilon_0Mc^3/(Q^2\omega_0^2)$ 即所求时间. 本稿校正: 官方(d)将能量损失率与振幅导数写为正值; 正确微分必须为负, 虽其最后时间数值仍相同. 官方辐射公式的三重叉积中间式也有多余负号, 本稿采用与最终矢量场一致的麦克斯韦符号.
\end{solution}
